% from ega2-1-fr.tex
\oldpage[II]{194}
Remplaçant $\mathcal{M}$ par $\widehat{\mathcal{M}}$, $\mathcal{S}$ par
$\widehat{\mathcal{S}}$ et $\mathcal{S}_X$ par
$(\mathcal{S}_X^{\geq})^{\wedge}$ dans le raisonnement précédent, on a de
même un isomorphisme canonique
\begin{equation}
\label{II.8.12.5.2-fr}
  q^*(\widehat{\mathcal{M}})
  \otimes_{q^*(\widehat{\mathcal{S}})}(\mathcal{S}_X^{\geq})^{\wedge}
  \xrightarrow{\sim}(\mathcal{M}_X^{\geq})^{\wedge}.
\tag{8.12.5.2}
\end{equation}
Si l'on se souvient (\hyperref[II.8.6.2-fr]{8.6.2}) que le morphisme
structural $u:\operatorname{Proj}(\mathcal{S}_X^{\geq})\to X$ est un
isomorphisme, on déduit d'abord de ce qui précède que l'on a un
$u$-isomorphisme canonique
\begin{equation}
\label{II.8.12.5.3-fr}
  \operatorname{Proj}_0\mathcal{M}
  \xrightarrow{\sim}\operatorname{Proj}_0(\mathcal{M}_X^{\geq})
\tag{8.12.5.3}
\end{equation}
comme cas particulier de (\hyperref[II.8.12.4.2-fr]{8.12.4.2}). On constate
en effet, avec les notations de la démonstration de
(\hyperref[II.8.6.2-fr]{8.6.2}), que cela revient à voir que
l'homomorphisme canonique
$M_{(f)}\otimes_{S_{(f)}}(S_f^{\geq})^{(d)}\to(M_f^{\geq})^{(d)}$ est un
isomorphisme lorsque $f\in S_d$, ce qui est immédiat.

En second lieu, l'isomorphisme (\hyperref[II.8.12.5.2-fr]{8.12.5.2}) permet,
en appliquant cette fois (\hyperref[II.8.12.4.3-fr]{8.12.4.3}) au morphisme
canonique $r=\operatorname{Proj}(\widehat{\alpha}):\widehat{C}_X\to
\widehat{C}$, d'obtenir un $r$-morphisme canonique
\begin{equation}
\label{II.8.12.5.4-fr}
  \mathcal{M}^{\square}\to(\mathcal{M}_X^{\geq})^{\square}.
\tag{8.12.5.4}
\end{equation}
Rappelons maintenant (\hyperref[II.8.6.2-fr]{8.6.2}) que les restrictions de
$r$ aux cônes épointés $\widehat{E}_X$ et $E_X$ sont des \emph{isomorphismes}
sur $\widehat{E}$ et $E$ respectivement. En outre~:
\end{env}

\begin{proposition}[8.12.6]
\label{II.8.12.6-fr}
Les restrictions à $\widehat{E}_X$ et à $E_X$ du $r$-morphisme canonique
(\hyperref[II.8.12.5.4-fr]{8.12.5.4}) sont des isomorphismes
\begin{equation}
\label{II.8.12.6.1-fr}
  \mathcal{M}^{\square}|\widehat{E}
  \xrightarrow{\sim}(\mathcal{M}_X^{\geq})^{\square}|\widehat{E}_X
\tag{8.12.6.1}
\end{equation}
\begin{equation}
\label{II.8.12.6.2-fr}
  \widetilde{\mathcal{M}}|E
  \xrightarrow{\sim}(\mathcal{M}_X^{\geq})^{\sim}|E_X.
\tag{8.12.6.2}
\end{equation}
\end{proposition}

\begin{proof}
On se ramène au cas où $Y$ est affine comme dans la démonstration de
(\hyperref[II.8.6.2-fr]{8.6.2})~; avec les notations de cette dernière, et
en se ramenant aux définitions (\hyperref[II.2.8.8-fr]{2.8.8}), on doit
montrer que l'homomorphisme canonique
\[
  \widehat{M}_{(f)}\otimes_{\widehat{S}_{(f)}}
  (S_f^{\geq})_{(f/1)}^{\wedge}
  \to(M\otimes_S S_f^{\geq})_{(f/1)}^{\wedge}
\]
est un isomorphisme~; mais en vertu de
(\hyperref[II.8.2.3.2-fr]{8.2.3.2}) et
(\hyperref[II.8.2.5.2-fr]{8.2.5.2}), le premier membre s'identifie
canoniquement à
$M_f^{\leq}\otimes_{S_f^{\leq}}(S_f^{\geq})_{f/1}^{\leq}$, donc à
$M_f^{\leq}$ en vertu de (\hyperref[II.8.2.7.2-fr]{8.2.7.2}), et le second à
$(M_f^{\geq})_{f/1}^{\leq}$, donc aussi à $M_f^{\leq}$ en vertu de
(\hyperref[II.8.2.9.2-fr]{8.2.9.2}), d'où la conclusion en ce qui concerne
(\hyperref[II.8.12.6.1-fr]{8.12.6.1})~; la formule
(\hyperref[II.8.12.6.2-fr]{8.12.6.2}) résulte alors de
(\hyperref[II.8.12.6.1-fr]{8.12.6.1}) et de
(\hyperref[II.8.12.2.1-fr]{8.12.2.1}).
\end{proof}

\begin{corollary}[8.12.7]
\label{II.8.12.7-fr}
Avec les identifications de (\hyperref[II.8.6.3-fr]{8.6.3}), la restriction
de $(\mathcal{M}_X^{\geq})^{\square}$ à $\widehat{E}_X$ s'identifie à
$(\mathcal{M}_X^{\leq})^{\sim}$ et la restriction de
$(\mathcal{M}_X^{\geq})^{\square}$ à $E_X$ s'identifie à
$\widetilde{\mathcal{M}}_X$.
\end{corollary}

\begin{proof}
On est en effet ramené au cas affine, et cela résulte de l'identification de
$(M_f^{\geq})_{f/1}^{\leq}$ à $M_f^{\leq}$ et de
$(M_f^{\geq})_{f/1}$ à $M_f$ (\hyperref[II.8.2.9.2-fr]{8.2.9.2}).
\end{proof}

\begin{proposition}[8.12.8]
\label{II.8.12.8-fr}
Sous les hypothèses de (\hyperref[II.8.6.4-fr]{8.6.4}), l'homomorphisme
canonique (\hyperref[II.8.12.3.1-fr]{8.12.3.1}) est un isomorphisme.
\end{proposition}

\begin{proof}
Compte tenu du fait que $\operatorname{Proj}(\mathcal{S}_X^{\geq})\to X$
est un isomorphisme (\hyperref[II.8.6.2-fr]{8.6.2}) et des

